
\section{Inclusion functor and the generators for the braid group}\label{sec: inc fun gens}



This section collects various facts about the matrix factorization convolution algebra, mostly property that do not require intense computations but these properties will be used later.
In particular, the relation (\ref{eq: commuting}) implies that matrix factorizations corresponding commuting elementary braids commute inside the convolution algebra.
% More generally, we show that two matrix factorizations corresponding to 
%braids  commute if they are tangle different sets of strands.





\subsection{Unit in convolution algebra}\label{sec: unit} In this section we describe the unit in the convolution algebra.
 Let us denote by $I_{--}\subset\CC[\calXr_2]$ the ideal generated by matrix coefficients of $(g)_{--}$ where $(X,g,Y)$ are coordinates on $\calXr_2$.
Since the generators of $I_{--}$ form a regular sequence
we can define $\bcalC_{\parallel}:=\mathrm{K}^{\Wr}(I_{--}).$ A direct computation shows that
$$\SymbolPrint{\calC_{\parallel}}:=\Phi_n(\bcalC_{\parallel})=\mathrm{K}^{W_1}(X_{--})\otimes \mathrm{K}^{W_2}(g_{--})\in \MF_{B_n^2}(\calX^{\circ}_2,W),$$
where $W_1=\Tr(X_{--}(Y_1-\Ad_g(Y_2)_{++})$, $W_2=\Tr(X(\Ad_g(Y_2)_{-})$.

\begin{proposition}\label{prop: unit} For any $\calF\in \MF_{B^2_n}(\calX^{\circ},W)$ we have
$$ \calC_{\parallel}\star \calF=\calF=\calF\star\calC_{\parallel}.$$
\end{proposition}
\begin{proof}
Let us prove the second equation, first is analogous.
Let us denote by $C$ the $B_n^3$-equivariant matrix factorization $\pi_{12}^{\circ*}(\calF)
\otimes \pi_{23}^{\circ *}(\calC_{\parallel})$. Let us fix the  coordinates $(X,g_{13},g_{23}, Y_1,Y_2,Y_3)$ on $\calX^{\circ}_3$. Since
$$C^\sharp=\mathrm{K}^{\pi_{23}^{\circ*}(W_1)}((g_{23})_{--})\otimes [\Ad_{g_{12}}(X)_{--},Y_2-\Ad_{g_{23}}(Y_3)]\otimes \pi_{12}^{\circ*}(\calF)$$
we can use the differential of the Koszul complex to eliminate variables $(g_{23})_{--}$. We use the row transformation of the first type so we do not need to worry about
admissibility.

 As the result of elimination we obtain matrix factorization $C_1$ which is the tensor product of two factors $C'_1\otimes C''_1$  where $C''_1=\pi_{12}^{\circ *}(\calF)|_{(g_{23})_{--}=0}$ and
 $C'_1$ has the following equivariant limit
 $$ (C'_1)^{\sharp}=[(g_{23})_{--},R_1]\otimes [Y_2-\Ad_{(g_{23})_+}(Y_3),R_2].$$
 Now the us notice that $\pi_{13}^{\circ *}(W)$ does not depends on $g_{23}$ and the only term in $C_1^\sharp$ that contains dependence on $(g_{23})_{--}$ is
 $[(g_{23})_{--},R_1]$ hence $R_1=0$.


 Let us introduce $\calX_{3}(B_n)=\frg_n\times G_n\times B_n\times \frn_n^3$.
 The variety $\calX_3(B_n)$ is isomorphic to the subvariety of $\calX_3$ defined by the equations $(g_{23})_{--}=0$.
 Now we can contract the complex $C_1$ along the differentials of the Koszul complex $[(g_{23})_{--},0]$ to show the homotopy equivalence of
 the $B_n^3$-equivariant  matrix factorizations over the ring $\pi_{13}^{\circ*}(\CC[\calX_2^{\circ}])$:
 $$C_1\sim C_2= [Y_2+\Ad_{g_{23}}(Y_3),R_2]\otimes \pi_{12}^{\circ*}(\calF)\in \MF_{B_n^3}(\calX_3(B_n),\pi_{13}^{\circ*}(W)),$$
 where the maps $\pi_{ij}^{\circ}$ are the restriction of the corresponding maps from $\calX_3$.

 As the next step we use the first factor in the matrix factorization $C_2$ to eliminate $Y_2$ from the rest of the factors. Since $Y_2+\Ad_{g_{23}}(Y_3)$ is $B_n^3$ equivariant we again only
 use the first type of row transformations. Similarly to the previous step, after elimination we obtain the matrix factorization
 $$ C_3=[Y_2+\Ad_{g_{23}}(Y_3),0]\otimes  \pi_{12}^{\circ*}(\calF)|_{Y_2=\Ad_{g_{23}}(Y_3)}.$$

 Now we can contract the complex along the differentials of the Koszul complex $[Y_2+\Ad_{g_{23}}(Y_3),0]$ to obtain the homotopy equivalence of
 of
 the$B_n^3$-equivariant  matrix factorizations over the ring $\pi_{13}^{\circ*}(\CC[\calX_2^{\circ}])$:
 \begin{equation*}
  C_3\sim C_4=\tilde{\pi}_{12}^{\circ^*}(\calF)\in \MF_{B_n^3}(\calX_2\times B_n,W).
  \end{equation*}
 where $\tilde{\pi}_{12}^{\circ}:\calX_2\times B_n\to \calX_2$ is given by $(X,g,Y_1,Y_2,b)\mapsto (X,Y_1,\Ad_{b}(Y_2))$.

 Thus we see that the differentials of $\tilde{\pi}_{12}^{\circ*}(\calF)$ are $B_n^{(2)}$ invariant and the computation of
 $\CE_{\frn_n^{(2)}}(C_4)$ reduces to the computation of
 $$\CE_{\frn_n^{(2)}}(\CC[\calX_2\times B_n])=\CC[\calX_2]\otimes \CE_{\frn_n^{(2)}}(\CC[B_n]).$$
 
Since \(Q_{k\dots n}=B_n\) for \(k=n\), the proposition~\ref{prop: Lie cohs} implies
 that $\CE_{\frn_n^{(2)}}(\CC[B_n])=\CC[T_n]$. Thus after taking $T_n^{(2)}$-invariant part of
 $\CE_{\frn_n^{(2)}}(C_4)$ we obtain the matrix factorization $\calF$.
\end{proof}

%\begin{proposition} \label{prop: unit} For every $i$ we have
%\begin{enumerate}
%\item $\bcalC^{(i)}_\parallel=\bcalC_\parallel$.
%\item $\bcalC_\parallel$ is the unit of the convolution algebra.
%\end{enumerate}
%\end{proposition}
%\begin{proof}
%The first part follows from our construction of the induction functor:
%$$\bcalC^{(i)}_\parallel:=i_{i,i+1}\circ p^*_{i,i+1}(\mathrm{K}^{\Wr_2}((g)_{21}))=
%\mathrm{K}^{p^*_{i,i+1}(\Wr_2)}(g_{i+1,i})\otimes \mathrm{K}^{\Wr-\Wr_2}(\Pi^{i,i+1}(g))=\bcalC_\parallel.$$

%Let us fix coordinates $(X,g_{13},g_{23},Y)$ on $\calXr_3$. The second part follows from
%\begin{multline*}
%\calF\bar{\star}\bcalC_{\parallel}=\brpi_{13*}(\CE_{\frn^{(2)}}(\brpi_{12}^*(\calF)\otimes \mathrm{K}^{\brpi_{23}^*(\Wr)}((g_{23})_{--})))^{T^{(2)}})=\\
%\brpi_{13*}(\CE_{\frn^{(2)}}(\brpi_{12}^*(\calF)|_{g_{23}\in B})^{T^{(2)}})=\brpi_{13*}(\brpi_{12}^*(\calF)|_{g_{23}\in B}\otimes [(g_{23})_{++},0])^{T^{(2)}}=\calF.
%\end{multline*}
%where the second equality is due to the fact we can eliminate the variables $(g_{23})_{--}$ from the rest of the tensor product using Koszul complex
%$ \mathrm{K}^{\brpi_{23}^*(\Wr)}((g_{23})_{--}))$ and the fact that $\brpi^*_{13}(\Wr)$ does not depend on $g_{23}$. The third equality is due  to proposition~\ref{prop: Lie cohs}.
%\end{proof}

\begin{corollary} The matrix factorization $\bcalC_{\parallel}$ is the unit in the convolution algebra $$(\MF_{B_n^2}(\calXr_2,\Wr),\bar{\star}).$$
\end{corollary}

Let us also provide an alternative proof of the proposition. The alternative argument is more geometric and uses  change of base. Let us first notice that
$$ \SymbolPrint{\bcalC_{\parallel}}=i_{*}^{\parallel}(\calO),$$
where $\SymbolPrint{i^{\parallel}}: \frg_n\times B\times \frn_n\to \calX_2$ is the embedding
$(X,b,Y)\to (X,b,Y,\Ad_b(Y))$ and $\calO\in \MF_{B_n^3}(\frb\times B\times \frn_n^2,0)$ is the trivial matrix
factorization.% Let us introduce notation $\calX_2^{\circ}(B):=\frb\times B\times \frn_n^2$ and $\calX_3^{\circ}(B):=\frb\times B\times G_n\times \frn_n^3$.
  Then we following diagram of maps:
  $$ \begin{tikzcd}
 \calX_2^{\circ}\times B \arrow[dotted]{r}{i^{\parallel}}\arrow[dotted]{d}{\pi_{12}^\circ\times\pi_{23}^\circ}   & \calX_3\arrow{d}{\pi_{12}^\circ\times\pi_{23}^\circ}\arrow{r}{\pi_{13}^{\circ}}&\calX^\circ_2\\
\calX_2^\circ\times \frg_n\times B\times \frn_n\arrow{r}{1\times i^{\parallel}} &                                  \calX^\circ_2\times\calX^\circ_2&
 \end{tikzcd}$$
 where the dotted maps are: $i^{\parallel}$ is the  inclusion  \[ i^{\parallel}(X,g,Y_1,Y_2)=(X,gb^{-1},b,Y_1,Y_2,\Ad_b(Y_2))\] and
 the restriction of the  map \(\pi_{12}^{\circ}\)
to the subspace $\calX_3(B)$: \[\pi_{12}^\circ\times \pi_{23}^\circ(X,g,Y_1,Y_2,b)=(X,g,Y_1,Y_2,X,b,Y_2).\]

Now let use base change property to simplify \( \calF\star\calC_\parallel\):
\begin{multline*}
\CE_{\frn_n^{(2)}}(\pi_{13*}^{\circ}((\pi_{12}^{\circ}\times\pi_{23})^*(\calF\boxtimes \calC_\parallel)))^{T^{(2)}}=
\CE_{\frn_n^{(2)}}(\pi_{13*}^{\circ}((\pi_{12}^{\circ}\times\pi_{23})^*(i_*^\parallel(\calF\boxtimes \calO))))^{T^{(2)}}\\=
\CE_{\frn_n^{(2)}}((\pi_{13}^{\circ}\circ i^\parallel)_*(\pi_{12}^{\circ}\times\pi_{23}^\circ)^*(\calF\boxtimes \calO))^{T^{(2)}}=\CE_{\frn^{(2)}}(\tilde{\pi}_{12}^{\circ*}(\calF))^{T^{(2)}}
\end{multline*}
The last step of this proof the same as in the previous proof.

\subsection{Inclusion functor}
In the proof of proposition~\ref{prop: ind trans} we constructed the functor $\ind_{k,m}$ and we can use this functor to define the functor:
$$\SymbolPrint{\Ind_{k\dots m}}: \MF_{ B_{m-k+1}^2}(\calX_2(G_{m-k+1}),W)\to \MF_{ B_n^2}(\calX_2(G_n),W),$$
by $\SymbolPrint{\Ind_{k\dots m}}(\calF):=\ind_{k,m}(\calC_{\parallel}\boxtimes\calF\boxtimes\calC_{\parallel})$.
Using Knorrer functor we also define the reduced version of the inclusion functor:
$$\SymbolPrint{\Indb_{k\dots m}}: \MF_{B_{m-k+1}^2}(\calXr_2(G_{m-k+1}),\Wr)\to \MF_{B_n^2}(\calXr_2(G_n),\Wr),$$


The results of the previous section  the properties of the unit immediately imply the following corollaries


\begin{corollary}\label{cor: ind} The functors $\Indb_{k\dots m}$ and $\Ind_{k\dots m}$ are homomorphism of the convolution algebras.
\end{corollary}

\begin{corollary}\label{cor: commute non red}For any $k,m,k',m'$ such that $k<m<k'<m'$ and
 \[\calF\in \MF_{B_{m-k+1}^2}(\calX_2(G_{m-k+1}),W),\]
\[\calG\in \MF_{ B_{m'-k'+1}^2}(\calX_2(G_{m'-k'+1}),W)\] we have
$$ \Ind_{k\dots m} (\calF)\star\Ind_{k'\dots m'} (\calG)=\Ind_{k'\dots m'} (\calG)\star\Ind_{k\dots m} (\calF).$$
\end{corollary}


\begin{corollary}\label{cor: commute}For any $k,m,k',m'$ such that $k<m<k'<m'$ and
 $$\calF\in \MF_{ B_{m-k+1}^2}(\calXr_2(G_{m-k+1}),\Wr),\quad
\calG\in \MF_{ B_{m'-k'+1}^2}(\calXr_2(G_{m'-k'+1}),\Wr)$$ we have
$$ \Indb_{k\dots m} (\calF)\bar{\star}\Indb_{k'\dots m'} (\calG)=\Indb_{k'\dots m'} (\calG)\bar{\star}\Indb_{k\dots m} (\calF).$$
\end{corollary}



\subsection{Generators of the braid group}\label{ssec:gens}

When $n=2$ the potential is $\Wr(X,g,Y)=\Tr(X\Ad_g(Y))$ 
where
$$ g=\begin{bmatrix} g_{11}&g_{12}\\ g_{21}& g_{22}\end{bmatrix},\quad X=\begin{bmatrix} x_{11}&x_{12}\\ 0& x_{22}\end{bmatrix},\quad Y=\begin{bmatrix} 0& y_{12}\\0&0\end{bmatrix}$$
A direct computation shows that it factors:
$$\Wr(X,g,Y)=y_{12}(g_{11}(x_{11}-x_{22})+g_{21}x_{12})g_{21}/\Delta,\quad \Delta=\det(g).$$

If \(n=2\) then the unit matrix factorization from section~\ref{sec: unit} is:
\[ \crX_{\parallel}=\Phi(\begin{bmatrix}g_{21}\Delta^{-1}& y_{12}(g_{11}(x_{11}-x_{22})+g_{21}x_{12})\end{bmatrix}),\]

Let also introduce the matrix factorizations $\crX_{+}$,  $\crX_{\bullet}$:
 \[\SymbolPrint{\crX_{\bullet}}=\Phi(\begin{bmatrix}y_{12}\Delta^{-1}& g_{21}(g_{11}(X_{11}-X_{22})+g_{21}X_{12})\end{bmatrix}),\]
\begin{equation}\label{eq: C+}
\SymbolPrint{\crX_{+}}=\Phi(\begin{bmatrix}g_{21}y_{12}\Delta^{-1}& (g_{11}(x_{11}-x_{22})+g_{21}x_{12})\end{bmatrix}).
\end{equation}
The complex $\crX_+$, as we will see later, correspond to the positive crossing of two strands in the braid.
To define the complex corresponding to the negative crossing we  need to shift the weights of $B^2$ action on $\crX_+$.

Let $\chi_1,\chi_2$ be the generators of the group of characters of $B$:
$$\SymbolPrint{\chi_1}\left(\begin{bmatrix}a&b\\0&c\end{bmatrix}\right)=a,\quad \SymbolPrint{\chi_2}\left(\begin{bmatrix}a&b\\0&c\end{bmatrix}\right)=c.$$
Bellow we use convention that $\calF\brwsh{\chi'}{\chi''}\in \MF_{B^2}(\calX_2(G_2),W)$  is obtained
from the complex $\calF$ by twisting the action of the first factor in $B^2$ by the character
$\chi'$ and the action of the second copy is twisted by the character $\chi''$. With these conventions in mind we have the
following formula for the complex of negative crossing.
\begin{equation}\label{eq: C-}
 \SymbolPrint{\crX_-}:=\crX_+\brwsh{-\chi_1}{\chi_2}.
\end{equation}

We postpone the description of the $T_{sc}$-equivariant structure of matrix factorizations from above till the
section~\ref{sec: two str}.
In the section~\ref{sec: two str} we show that the convolution of  $\crX_{+}$  and $\crX_-$  is homotopic to \(\crX_\parallel\).

To define the generators of $\Brgr_n$ for general $n$ we use our induction functor:
$$\SymbolPrint{\crX_{\pm}^{(i)}}:=\Ind_{ii+1}(\crX_{\pm}),\quad \crX_{\parallel}^{(i)}:=\Ind_{ii+1}(\crX_\parallel).$$


 It is elementary to check that $\crX_{\parallel}^{(i)}$ is independent of $i$ and is a unit in the convolution algebra (see next section for the proof).
The main result of the first part of the paper is the following theorem

\begin{theorem}\label{thm: braids} The complexes $\crX_{\pm}^{(i)}$ satisfy the braid relations
\begin{gather}
\crX_{+}^{(i)}\star\crX_{-}^{(i)}\sim\crX_\parallel,\label{eq: inverse}\\
\crX_{+}^{(i)}\star \crX_{+}^{(i+1)} \star \crX_{+}^{(i)}\sim \crX_{+}^{(i+1)}\star \crX_{+}^{(i)}\star \crX_{+}^{(i+1)},\label{eq: cubic}\\
\crX_{+}^{(i)}\star \crX_{+}^{(j)}\sim\crX_{+}^{(j)}\star \crX_{+}^{(i)},\quad |i-j|>1.\label{eq: commuting}
\end{gather}
where \(\sim\) is the homotopy equivalence.
\end{theorem}
\begin{proof}
The relation (\ref{eq: inverse}) is proven for two strand braids in lemma~\ref{thm: two strands}. Hence by corollary~\ref{cor: ind} the relation (\ref{eq: inverse}) holds in general.
The equation (\ref{eq: cubic}) is proven for three strand braids in corollary~\ref{cor: cubic}. Finally, the commuting relation (\ref{eq: commuting}) follows from
the corollary~\ref{cor: commute}.
\end{proof}

For a given element $\beta=\sigma_{i_1}^{\epsilon_1}\cdots\sigma_{i_\ell}^{\epsilon_\ell}$ we introduce the following notations:
$$ \calC_\beta=\calC_{\epsilon_1}^{(i_1)}\star\dots \calC_{\epsilon_\ell}^{(i_\ell)},\quad \bcalC_\beta=\bcalC_{\epsilon_1}^{(i_1)}\bar{\star}\dots \bar{\star}\bcalC_{\epsilon_\ell}^{(i_\ell)}.$$





\section{General properties of convolution algebra}\label{sec: aux sec}


\subsection{Auxiliary Lemma} In the proof of the braid relations and Markov moves we will use some generalization of the previous corollary which we explain below.
The statement that we prove in this section allows us to use smaller intermediate space for computation of the convolution.

Let us introduce notation $\SymbolPrint{Q_{k\dots m}}\subset G_n$ for subgroup with Lie algebra $\SymbolPrint{\frq_{k\dots m}}$ spanned by $E_{ij}$, $i\le j$, $i,j\in [1,n]$ and by $E_{ij}$, $i>j$,
$i,j\in [k,m]$. Respectively, we denote by $\SymbolPrint{\Pi_+^{\overline{k\dots m}}}$ the projection map  $\Lie(G_n)\to \frq_{k\dots m}$ and $\SymbolPrint{\Pi_{--}^{\overline{k\dots m}}}$ is defined by
$\mathrm{Id}=\Pi_+^{\overline{k\dots m}}+\Pi_{--}^{\overline{k\dots m}}$. Similarly we define $\SymbolPrint{\Pi_{++}^{\overline{k\dots m}}}(X)=\Pi_{--}^{\overline{k\dots m}}(X^t)^t$, where $t$ stands for transpose.  Let us also denote by $\SymbolPrint{\Pi^{\overline{k\dots m}}}$ the projection $\frg_n\to \frg_{m-k+1}$ that extracts the block of the matrix elements
with entries $ij$, $i,j\in [k,m]$. We use the same notation for the maps of the corresponding groups.





Let us introduce the space
$\SymbolPrint{\calX^\circ_3(G_n,G_{k\dots m})}=\frg_n\times G_n\times G_{k-m+1}\times \frn_n\times \frn_{k-m+1}\times \frn_n$. Let us also fix an embedding
$\SymbolPrint{i_{k\dots m}}$ of this space inside $\calX_3$:

\[ i_{k\dots m}(X,g_{12},g_{23},\vec{Y})=(X,g_{12},i_{k\dots m}(g_{23}),Y_1,i_{k\dots m}(Y_2)+\Pi_{++}^{\overline{k\dots m}}(\Ad_{g_{23}}(Y_3)),Y_3),\]
where \(\vec{Y}=(Y_1,Y_2,Y_3)\), \(i_{k\dots m}:G_{k-m+1}\to G_n\) and \(i_{k\dots m}: \frn_{k-m+1}\to \frn_n\) are the natural embeddings.

Respectively we define maps \[\pi_{12}^{\circ}:\calX^\circ_3(G_n,G_{k\dots m})\to \calX_2(G_{n}),\]
\[\pi_{23}^{\circ}:\calX^\circ_3(G_n,G_{k\dots m})\to \calX_2(G_{m-k+1}),\] \[\pi_{13}^{\circ}:\calX^\circ_3(G_n,G_{k\dots m})\to \calX_2(G_n),\]
\begin{equation}\label{eq: pi12}
\pi_{12}^{\circ}(X,g_{12},g_{23},Y_1,Y_2,Y_3)=(X,g_{12},Y_1,i_{k\dots m}(Y_2)+\Pi_{++}^{\overline{k\dots m}}(\Ad_{g_{23}}(Y_3)).
\end{equation}
\begin{equation}\label{eq: pi23}
\pi_{23}^{\circ}(X,g_{12},g_{23},Y_1,Y_2,Y_3)=(\Pi^{\overline{k\dots m}}( \Ad_{g_{12}}^{-1}(X)),g_{23},Y_2,\Pi^{\overline{k\dots m}}(Y_3)),
\end{equation}
$$\pi_{13}^{\circ}(X,g_{12},g_{23},Y_1,Y_2,Y_3)=(X,g_{12}g_{23},Y_1,Y_3).$$
These maps are the restrictions of the corresponding maps from $\calX_3$ on the embedded subvariety $\calX^\circ_3(G_n,G_{k\dots m})$.

Similarly, we define $\calX^\circ_3(G_{k\dots m},G_n)=\frg_n\times G_{k\dots m}\times G_n\times \frn_n\times \frn_{k-m+1}\times \frn_n$. Let us also fix an embedding $i_{k\dots m}$ of this space inside $\calX_3$:
$$ i_{k\dots m}(X,g_{12},g_{23},\vec{Y})=(X,i_{k\dots m}(g_{12}),g_{23},Y_1,i_{k\dots m}(Y_2)+\Pi_{++}^{\overline{k\dots m}}(\Ad^{-1}_{g_{12}}(Y_1)),Y_3),$$
\(\vec{Y}=(Y_1,Y_2,Y_3)\).
The maps $\pi_{ij}^{\circ}$ are defined analogously to the previous case: by the restriction from ambient space.

The spaces $\calX^\circ_3(G_n,G_{k\dots m})$ and $\SymbolPrint{\calX^\circ_3(G_{k\dots m},G_n)}$ have unique $B_n^3$ equivariant structure that makes the maps $\pi_{ij}^\circ$ $B_n^3$-equivariant.

\begin{lemma}\label{prop: small flag}  For any \[\calF\in \MF_{B_{m-k+1}^2}(\calX_2(G_{m-k+1}),W),\calG\in \MF_{B_n^2}(\calX_2(G_n),W)\] we have
$$ \calG\star\Ind_{k\dots m}(\calF)=\CE_{\frn_{m-k+1}^{(2)}}\pi_{13*}^\circ(\pi_{12}^{\circ*}(\calG)\otimes \pi_{23}^{\circ*}(\calF))^{T^{(2)}_{m-k+1}},$$
$$ \Ind_{k\dots m}(\calF)\star \calG=\CE_{\frn_{m-k+1}^{(2)}}\pi_{13*}^\circ(\pi_{12}^{\circ*}(\calF)\otimes \pi_{23}^{\circ*}(\calG))^{T^{(2)}_{m-k+1}},$$
\end{lemma}
The key observation in the proof is an alternative description of the induction functor $\Ind_{k\dots m}$. The space $\SymbolPrint{\calX^\bullet_2(Q_{k\dots m})}:=
\frg_n\times Q_{k\dots m}\times \frn_{m-k+1}\times \frn_n$ can be mapped into
$\calX_2(G_n)$ with the map:
\[ \SymbolPrint{j_{k\dots m}}(X,g,Y_1,Y_2)=(X,j_{k\dots m}(g),i_{k\dots m}(Y_1)+\Pi_{++}^{\overline{k\dots m}}(\Ad_g(Y_2)),Y_2)\]
where \(j_{k\dots m}: Q_{k\dots m}\to G_n\) is the natural embedding.

There is also a natural projection map $\Pi^{\overline{k\dots m}}:\calX^\bullet_2(Q_{k\dots m})\to \calX_2^\circ(G_{m-k+1})$ defined by:
$$ \Pi^{\overline{k\dots m}}(X,g,Y_1,Y_2)=(\Pi^{\overline{k\dots m}}(X),\Pi^{\overline{k\dots m}}(g),\Pi^{\overline{k\dots m}}(Y_1),\Pi^{\overline{k\dots m}}(Y_2)).$$
Both of these maps are $B_n^3$-equivariant. In particular, we have a commuting diagram:
  \[ \begin{tikzcd}
 \calX_2^{\bullet}(Q_{k\dots m}) \arrow[dotted]{r}{i'}\arrow[dotted]{d}{\Pi'_1\times\Pi^{\overline{k\dots m}}\times\Pi'_3}   & \calX^\circ_2(P_{k,m})\arrow{d}{p_{k,m}}\arrow{r}{i_{k,m}}&\calX''_{2,n}\\
 \calX'_{2,k}\times \calX''_{2,m-k+1}\times \calX'_{2,n-m} \arrow{r}{i^{\parallel}\times 1\times i^{\parallel}}\arrow{d}{\pi_2} &
 \calX''_{2,k}\times\calX''_{2,m-k+1}\times \calX''_{2,n-m}&\\
\calX^\circ_2(G_{m-k+1})&&
\end{tikzcd}\]
where \(\calX'_{2,k}:=\frg_k\times B_k\times \frb_k\times \frb_k=\calX^{\circ}_2(B_k)\), \(\calX''_{2,k}:=\frg_k\times G_k\times \frb_k\times\frb_k=\calX^{\circ}_2(G_k)\)
and the dotted arrow maps are:
\[i'(X,g,Y_1,Y_2)=(\Pi^{\overline{k\dots m}}_+(X),j_{k\dots m}(g),i_{k\dots m}(Y_1)+
  \Pi^{\overline{k\dots m}}_{++}(\Ad_g(Y_2)),Y_2),\]
\[\Pi'_i(X,g,Y_1,Y_2)=(p_{k,m}^{(i)}(X),p_{k,m}^{(i)}(g),p_{k,m}^{(i))}(\Ad_g(Y_2)),p_{k,m}^{(i)}(Y_2)),\]
here \(p_{k,m}^{(1)}\), \(p_{k,m}^{(3)}\) are the compositions of \(p_{k,m}\) with the projections on \(k\times k\) and \((n-m)\times(n-m)\) blocks.

% Then our uniqueness theorem for the push forward implies:
Using the base change in our previous diagram we obtain
\begin{multline*}
\Ind_{k\dots m}=i_{k,m,*}\circ p_{k,m}^*\circ (i^{\parallel}\times 1\times i^{\parallel})\circ \pi_2^*=i_{k,m,*}\circ i'_*\circ (\Pi'_1\times\Pi^{\overline{k\dots m}}\times\Pi'_3)^*\circ\pi_2^*\\=
j_{k\dots m,*}\circ \Pi^{\overline{k\dots m},*}.
\end{multline*}

\begin{proof}[Proof of proposition~\ref{prop: small flag}]
We prove the first formula the second is analogous.
We have the following commuting diagram of the maps:
  $$ \begin{tikzcd}
 \calX_3^{\circ}(G_n,Q_{k\dots m}) \arrow[dotted]{r}{i_{k\dots m}}\arrow[dotted]{d}{\pi_{12}^\circ\times\pi_{23}^\circ}   & \calX^\circ_3\arrow{d}{\pi_{12}^\circ\times\pi_{23}^\circ}\arrow{r}{\pi_{13}^{\circ}}&\calX^\circ_2\\
\calX_2^\circ\times \calX_2^\bullet(Q_{k\dots m}) \arrow{r}{1\times j_{k\dots m}}\arrow{d}{1\times\Pi^{\overline{k\dots m}}} &                                  \calX^\circ_2\times\calX^\circ_2&\\
\calX^\circ_2(G_n)\times\calX^\circ_2(G_{m-k+1})&&
 \end{tikzcd}$$
where the dotted arrow map $\pi_{12}^\circ\times\pi_{23}^\circ$ is given by
$$ \pi_{12}^{\circ}(X,g_{12},g_{23},Y_1,Y_2,Y_3)=(X,g_{12},Y_1,i_{k\dots m}(Y_2)+\Pi_{++}^{\overline{k\dots m}}(\Ad_{g_{23}}(Y_3)),$$
$$\pi_{23}^{\circ}(X,g_{12},g_{23},Y_1,Y_2,Y_3)=(\Ad_{g_{12}}^{-1}(X),g_{23},Y_2,\Pi_{-}^{\overline{k\dots m}}(Y_3)).$$

Now we can use the base change formula to simplify \( \calG\star\Ind_{k\dots m}(\calF)\) :
\begin{multline*}
\CE_{\frn_n^{(2)}}(\pi^\circ_{13*}\circ (\pi^\circ_{12}\times \pi_{23}^\circ)^*\circ (\mathrm{Id}\times j_{k\dots m})_*\circ (\mathrm{Id}\times\Pi^{\overline{k\dots m}})^*(\calG\boxtimes \calF)))^{T^{(2)}_n}=\\
 \CE_{\frn_n^{(2)}}(\pi^\circ_{13*} \circ  (i_{k\dots m})_*\circ  (\pi^\circ_{12}\times \pi_{23}^\circ)^* \circ(\mathrm{Id}\times\Pi^{\overline{k\dots m}})^*(\calG\boxtimes \calF)))^{T^{(2)}_n}=\\
  \CE_{\frn_n^{(2)}}(\pi^\circ_{13*}\circ (\pi^\circ_{12}\times \pi_{23}^\circ)^*(\calG\boxtimes \calF)))^{T^{(2)}_n},
\end{multline*}
where in the last equation we use $ \pi^\circ_{13}\circ i_{k\dots m}=\pi_{13}^\circ$ and
 $\mathrm{Id}\times\Pi^{\overline{k\dots m}}\circ \pi^\circ_{12}\times \pi_{23}^\circ=\pi^\circ_{12}\times \pi_{23}^\circ$ where the maps in the product
$\pi^\circ_{12}\times \pi_{23}^\circ: \calX_3^\circ(G_n,Q_{k\dots m})\to \calX^\circ_2(G_n)\times\calX^\circ_2(G_{m-k+1})$ are given by the formulas (\ref{eq: pi12}), (\ref{eq: pi23}).

Let us denote by $C$ the matrix factorization $\pi^\circ_{13*}\circ (\pi^\circ_{12}\times \pi_{23}^\circ)^*(\calG\boxtimes \calF)$
For last step  we  use Hochschild-Serre spectral sequence to compute functor $\CE_{\frn_n^{(2)}}$:
$$ \CE_{\frn_{m-k+1}^{(2)}}(\CE_{\Lie(J'_{k\dots m})}(C))\Rightarrow \CE_{\frn_n^{(2)}}(C) $$
where $J_{k\dots m}$ is the kernel of the projection homomorphism $Q_{k\dots m}\to G_{m-k+1}$ and $J'_{k\dots m}$ is its unipotent part.  Let us abbreviate the quotient $J_{k\dots m}/J'_{k\dots m}$ by
$J/J'$.
We have
\begin{equation}\label{eq: CE}
 \CE_{\Lie(J'_{k\dots m})}(C)^{J/J'}=\pi^\circ_{13*}(\CE_{\Lie(J'_{k\dots m})}(\pi^\circ_{12}\times \pi_{23}^\circ)^*(\calG\boxtimes \calF)).
\end{equation}
On the other hand the formulas for the maps $\pi^\circ_{12}$ and $\pi^\circ_{23}$ imply that the differentials of the pull back $(\pi^\circ_{12}\times \pi_{23}^\circ)^*(\calG\boxtimes \calF)$
are invariant with respect to the action of $J'_{k\dots m}$. Hence to complete the computation of (\ref{eq: CE}) we need to  compute
\[ \mathrm{H}_{\Lie}^*(J'_{k\dots m},\CC[\calX^\circ_3(G_n,Q_{k\dots m})])^{J/J'}\] which is a tensor product of \(\CC[\frb\times G_n\times \frn_{m-k+1}\times \frn_n]\) and \(\mathrm{H}_{\Lie}^*(J'_{k\dots m},\CC[Q_{k\dots m}])^{J/J'},\)
and in proposition below we show that  $\CE_{J'_{k\dots m}}(\CC[Q_{k\dots m}])^{J'/J}=\CC[G_{m-k+1}].$
Hence the statement follows since $T^{(2)}_n=(J/J')\times T_{m-k+1}$ where $T_{m-k+1}$ is the maximal torus of $G_{m-k+1}\subset G_n$.



\end{proof}



\begin{proposition}\label{prop: Lie cohs} For any $k,m$, $1\le k<m\le n$ we have
$$\CE_{\Lie(J'_{k\dots m})}(\CC[Q_{k\dots m}])=\CC[G_{m-k+1}]\otimes \CC[J/J'].$$
\end{proposition}
\begin{proof}
We show slightly more general statement
$$\CE_{\Lie(J'_{k\dots m})}(\CC[Q_{k'\dots m'}])=
\CC[Q_{k-k'+1\dots m-k'+1}\subset G_{m-k+1}]\otimes \CC[(J/J')],$$
where $k\ge k'$, $m'\le m$.

We use induction by $m-k$ and the Hochschild-Serre Spectral sequence (\cite{Wei}, sect 7.5). Since we have short exact sequence of groups
$$ 1\to J'_{k+1,\dots, m}\longrightarrow J'_{k\dots n}\longrightarrow \CC^{m-k+1}\to 1,$$
the spectra sequence tells us that
$$\CE_{\CC^{m-k+1}}(\CE_{\Lie\left(J'_{k+1\dots m}\right)}(\CC[Q_{\k'\dots m'}]))\Rightarrow \CE_{J_{k\dots m}}(\CC[Q_{k'\dots m'}]).$$
By induction we have
$\CE_{\Lie(J'_{k+1\dots m})}(\CC[Q_{k'\dots m'}])=\CC[Q_{k-k'+2\dots m-k'+1}\subset G_{m-k+2}]$.
After relabeling matrix indices our statement is equivalent to the computation of the  following cohomology
$\mathrm{H}^*_{\Lie}(\CC^{m-k+1},\CC[Q_{k-k'+2\dots m-k'+1}])$. The ring
of polynomial functions on $Q_{k-k'+2\dots m-k'+1}$ is generated by the
elements $E_{kl}^\vee$    %\in \Hom(Q_{k-k'+2\dots m-k'+1},\CC)$
dual to the matrix units $E_{kl}$ of $E$. The generators
$\partial_i$ of the commutative algebra $\CC^{m-k+1}$ act on the ring generators by the formula:
$$\partial_i E^\vee_{kl}=\delta_{k1} E_{il}^\vee.$$
We can apply invertible linear transformation to our system of differential operators to obtain new system of generators:
$$\tilde{\partial}_k=\sum_{i=2}^{m-k+2} (-1)^{i+k}\Delta_{ik}\partial_k,$$
where $\Delta_{ik}\in \CC[Q_{k-k'+2\dots m-k'+1}]$ is the determinant of $(m-k)\times (m-k)$ obtained from $E$ by removing columns and rows containing $11$ and $ik$ entries.
Then  we have:
$$\tilde{\partial}_i E^\vee_{kl}=\delta_{k1}\delta_{il} E_{11}^\vee.$$
and by Poincare lemma we obtain
$$\mathrm{H}^*_{\Lie}(\CC^{m-k+1},\CC[Q_{k-k'+2\dots m-k'+1}])=\CC[Q_{k-k'+2\dots m-k'+1}]/(E_{12}^\vee,\dots,E_{1,m-k+2}^\vee)$$ and the
statement follows.
\end{proof}


\subsection{Reduced case version of the auxiliary lemma }

%Let us use notation $G_{k\dots m}\subset \mathbf{P}_{k\dots m}\cap G_n$ for the subgroup of elements $g$ such that $g_{ij}=0$ if $i\ne j$ and either $ i\notin  [k,m]$ or $j\notin [k,m]$.
%As an algebraic group it is isomorphic to $G_{m-k+1}\times(\CC^*)^{n-m+k-2}$.

Let us introduce \[\SymbolPrint{\calXr^{\circ}_2(G_{k\dots m})}=\frb\times G_{k\dots m} \times \frn^2,\]
\[\SymbolPrint{\calXr_3(G_n,G_{k\dots m})}:=\frb_n\times G_n\times G_{m-k+1}\times \frn_n.\] Let us also define an embedding
$\SymbolPrint{\bar{i}_{k\dots m}}:\calXr_3(G_n,G_{k\dots m})\to \calXr_3$ by
$$ \bar{i}_{k\dots m}(X,g_{12},g_{23},Y)=(X,g_{12},i_{k\dots m}(g_{23}),Y).$$  Using the last embedding we define the maps $\SymbolPrint{\bar{\pi}_{ij}}$ by restricting from the ambient $\calX_3$:
$$ \bar{\pi}_{12},\bar{\pi}_{13}: \calXr_3(G_n,G_{k\dots m})\to \calXr_2(G_n),\quad \bar{\pi}_{23}: \calXr_3(G_n,G_{k\dots m})\to \calXr_2(G_{k\dots m}).$$

Similarly, we define $\calXr_3(G_{k\dots m},G_n)$ together with the embedding $\bar{i}_{k\dots m}: \calXr_3(G_{k\dots m}, G_n)\to \calX_3$:
$$ \bar{i}_{k\dots m}(X,g_{12},g_{23},Y)=(X,i_{k\dots m}(g_{12}),g_{23},Y).$$
The same way as before we define the corresponding maps $\bar{\pi}_{ij}$.

The spaces $\calXr_3(G_{k\dots m},G_n)$ and $\calXr_3(G_n,G_{k\dots m})$ have natural $B_n\times B_{m-k+1}\times B_n$-equivariant structure. Moreover, the maps $\bar{\pi}_{ij}$ are $B_n^2$-equivariant but not $B_{m-k+1}$-equivariant.
However, we can introduce the $B_n\times B_{m-k+1}\times B_n$-equivariant structure on the pull-back $(\bar{\pi}_{12}\times \bar{\pi}_{23})(\calF\otimes\calG)$ where $\calF,\calG$ are equivariant matrix factorizations
from the appropriate spaces. More precisely, repeating the argument of the proposition~\ref{prop: Knorrer factor} we obtain

\begin{proposition}\label{prop: small prod 1}  Let
  \[\calF\in \MF_{B_n^2}(\calXr_2(G_n),\Wr), \quad\calG\in \MF_{B_{m-k+1}^2}(\calXr_2(G_{m-k+1}),\Wr)\] then there is a unique matrix factorization
$$(\bar{\pi}_{12}\otimes_{B_{m-k+1}}\bar{\pi}_{23})^*(\calF\boxtimes \calG)\in \MF_{B_n\times B_{m-k+1}\times B_n}(\calXr_3(G_n,G_{m\dots k}),\bar{\pi}_{13}^*(\Wr)),$$
such that its $B_n^2$-equivariant specialization  has properties:
$$(\bar{\pi}_{12}\otimes_{B_{m-k+1}}\bar{\pi}_{23})^*(\calF\boxtimes \calG)^\sharp=(\bar{\pi}_{12}\times\bar{\pi}_{23})^*(\calF\boxtimes \calG),$$
and
$$\CE_{\frn_{m-k+1}}(\pi_{13*}((\bar{\pi}_{12}\otimes_{B_{m-k+1}}\bar{\pi}_{23})^*(\calF\boxtimes \calG)))^{T_{m-k+1}^{(2)}}=\calF\bar{\star}\Indb_{k\dots m}(\calG).$$
\end{proposition}


\begin{lemma}\label{prop: small prod 2}  Let \[\calF\in \MF_{B_{m-k+1}^2}(\calXr_2(G_{m-k+1}),\Wr),\quad\calG\in \MF_{B_{m-k+1}^2}(\calXr_2(G_n),\Wr)\] then there is a unique matrix factorization
$$(\bar{\pi}_{12}\otimes_{B_{m-k+1}}\bar{\pi}_{23})^*(\calF\boxtimes \calG)\in \MF_{B_n\times B_{m-k+1}\times B_n}(\calXr_3(G_{m\dots k},G_n),\bar{\pi}_{13}^*(\Wr)),$$
such that its $B_n^2$ limit has properties:
$$(\bar{\pi}_{12}\otimes_{B_{m-k+1}}\bar{\pi}_{23})^*(\calF\boxtimes \calG)^\sharp=(\bar{\pi}_{12}\times\bar{\pi}_{23})^*(\calF\boxtimes \calG),$$
and
$$\CE_{\frn_{m-k+1}}(\pi_{13*}((\bar{\pi}_{12}\otimes_{B_{m-k+1}}\bar{\pi}_{23})^*(\calF\boxtimes \calG)))^{T_{m-k+1}^{(2)}}=\Indb_{k\dots m}(\calF)\bar{\star}\calG.$$
\end{lemma}


\subsection{Reduced convolution with rank $1$ matrix factorizations} \label{ssec: conv rk 1}
The most important case of these theorems is the case $m=k+1$, that is when one of the matrix factorizations in
$\bar{\star}$ product is lifted from the space $\calXr_2(G_2)$, we call such matrix factorization rank $1$ matrix factorization. In this case the pull back functor $(\bar{\pi}_{12}\times_{B_{m-k+1}}\bar{\pi}_{23})^*$ could be made very explicit.
We do not provide the details of the computation, we essentially just follow the method of proposition~\ref{prop: Knorrer factor}.

Let us fix notations. Let $\calF=(M_1,D_1,\partial'_l,\partial'_r)\in \MF_{B_n\times B_2}(\calXr_2(G_n),\Wr)$, $\calG=(M_2,D_2,\partial''_l,\partial''_r)\in \MF_{B_2\times B_n}(\calXr_2(G_2),\Wr)$, $B_2$ in both cases
is the image of the embedding $i_{k,k+1}: B_2\to B_n$.  Since the Lie algebra $\frn_2$ is one-dimensional  $\frn_2=\langle e \rangle$ the corresponding $\frn_2$-equivariant structure is given by the
maps the even maps $$\partial'_r\in Hom_{\CC[\calXr_2(G_n)]}(M_1, M_1), \partial''_l\in Hom_{\CC[\calXr_2(G_2)]}(M_2, M_2).$$ Respectively
\begin{equation}\label{eq: pi12Bpi23 1}
 (\bar{\pi}_{12}\otimes_{B_{m-k+1}}\bar{\pi}_{23})^*(\calF\boxtimes \calG)=(\bar{\pi}_{12}^*(M_1)\otimes\bar{\pi}_{23}^*(M_2),D_1+D_2;\partial'_l,\partial'_r+\partial''_l+\partial_{rl},\partial'_r),
 \end{equation}
where the differential $\partial_{rl}$ is given by the formula
\begin{equation}\label{eq: pi12Bpi23 2}
\partial_{rl}=\bar{\pi}_{12}^*\left(\frac{\partial D_1}{\partial Y_{k,k+1}}\right)(\tilde{Y}^2_{k+1,k+1}-\tilde{Y}^2_{k,k})+ \bar{\pi}_{23}^*\left(\frac{\partial D_2}{\partial X_{kk}}-\frac{\partial D_2}{\partial X_{k+1,k+1}}\right) \tilde{X}^2_{k+1,k},
\end{equation}
where $\tilde{X}^2=\Ad_{g_{12}}^{-1}(X)$ and $\tilde{Y}^2=\Ad_{g_{12}}(Y)$.

In the other case  \[\calF=(M_1,D_1,\partial'_l,\partial'_r)\in \MF_{B_n\times B_2}(\calXr_2(G_2),\Wr),\]
\[ \calG=(M_2,D_2,\partial''_l,\partial''_r)\in \MF_{B_2\times B_n}(\calXr_2(G_n),\Wr)\] and the expressions for
$(\bar{\pi}_{12}\times_{B_{m-k+1}}\bar{\pi}_{23})^*(\calF\boxtimes \calG)$ is given by the same formulas (\ref{eq: pi12Bpi23 1}), (\ref{eq: pi12Bpi23 2}).

%The $B_n^2$-equvariant structure on $\calXr_3$ restricts to $B_n\times B_{m-k+1}$ equivariant structure on  $\calXr_3(G_n,P_{k\dots m})$.

%\begin{lemma} Let $\calF\in \MF_{B^2_{m-k+1}}(\calXr_2(G_{m-k+1}),W)$, $\calG\in  \MF_{B_n^2}(\calXr_2(G_n),W)$ then
%$$\calG\bar{\star}\Indb_{k\dots m}(\calF)=\bar{\pi}_{13*}(\CE_{\frn^{(2)}_{m-k+1}}(\bar{\pi}_{12}^*(\calG)\otimes\bar{\pi}_{23}^*(\calF))^{T^{(2)}})$$
%\end{lemma}

\begin{remark} As one can see, if we work with $B_2$-equivariant matrix factorizations on $Z$ with $B_2$-invariant potential then all the technicalities that are discussed in the first few sections are
unnecessary since $\MFs_{B_2}^{str}(Z,F)=\MF_{B_2}(Z,F)$.
\end{remark}

\subsection{Involution} Let $w_0\in G_n$ be the longest element of the Coxeter group $S_n$: $(w_0)_{i,j}=\delta_{i+j,n+1}$.   The involution $\SymbolPrint{\mathrm{D}}: G\to G$ defined by
$\mathrm{D}(g)=w_0 (g^{-1})^t w_0$ preserves $B$ and the Lie version of involution $\mathrm{D}(X)=-w_0 X^t w_0$, \(X\in \frg=\textup{Lie}(G)\) preserves $\frb\subset \frg$. Moreover, since $Tr(\mathrm{D}(X))=-Tr(X)$ the
have automorphism of category \[\mathrm{D}:\quad \MF_{B^2}(\calX_2(G),W)\to \MF_{B^2}(\calX_2(G),W),\]
\[ \mathrm{D}:\quad \MF_{B^2}(\calXr_2(G),\Wr)\to \MF_{B^2}(\calXr_2(G),\Wr).\]
The following proposition is almost tautological

\begin{proposition} Automorphism $\mathrm{D}$ respects convolution product:
$$\mathrm{D}(\calF)\star \mathrm{D}(\calG)=\mathrm{D}(\calF\star\calG),\quad \mathrm{D}(\calF)\bar{\star} \mathrm{D}(\calG)=\mathrm{D}(\calF\bar{\star}\calG),$$
and
$$ \mathrm{D}(\calC^{(i)}_{\pm})=\calC^{(n-i)}_{\pm},\quad \mathrm{D}(\bcalC^{(i)}_{\pm})=\bcalC^{(n-i)}_{\pm}.$$
\end{proposition}
\begin{proof}
  Let us show that \(\mathrm{D}\) induces an automorphism of the category of matrix factorizations. We define the action on the objects
  as
  \[\calF=(M,\tilde{D},\partial_l,\partial_r)\mapsto (M,\mathrm{D}\circ \tilde{D}\circ \mathrm{ D},
    \mathrm{D}\circ \partial_l\circ \mathrm{D},\mathrm{D}\circ\partial_r\circ \mathrm{D}).\]
  Since \(W\circ \mathrm{D}=W\) the last formula sends objects of \(\MF_{B^2}(\calX_2(G),W)\) to the objects of \(\MF_{B^2}(\calX_2(G),W\circ\mathrm{D})=
  \MF_{B^2}(\calX_2(G),W)\). The action on the space of morphisms  is also defined by conjugation with \(\mathrm{D}\):
  \[ Hom(\calF,\calF') \owns \Psi\mapsto \mathrm{D}\circ \Psi\circ \mathrm{D}.\]
  Thus is constructed an automorphism of the category.

  The induction functors \(\Ind_{k,m},\ind_{k,m}\) are defined with pull-backs and push-forwards along the maps \(p_{k,m},i_{k,m}\)
  Hence the rest of the proposition follows from the observation that \(\mathrm{D}\circ i_{k,m}= i_{k',m'}\circ \mathrm{D}\), \(
  \mathrm{D}\circ p_{k,m}= p_{k',m'}\circ \mathrm{D}\),
  \(k'=n+1-k,m'=n+1-m\).
\end{proof}

